\section{Síntesis y ampliación}

\subsection{Conclusiones principales}

\begin{enumerate}
\item En 2025 la innovación en arquitectura vuelve a ser importante porque el crecimiento de los datos se desacelera, aumenta la demanda de contexto largo y el costo de inferencia se vuelve más pesado.
\item Full Attention tiene una gran capacidad expresiva, pero en secuencias largas tanto el cómputo como la KV cache resultan costosos.
\item La idea central de Linear Attention es reescribir la atención cuadrática como una recurrence casi lineal o como una actualización de estado.
\item La KDA de Kimi Linear usa un decay de grano más fino para aprovechar mejor la memoria de un hidden state limitado.
\item Kimi opta por Hybrid Linear, DeepSeek por Sparse y MiniMax M2 regresa a Full, lo que muestra que Attention todavía no ha convergido.
\item La FFN ya se transformó con MoE; Attention es la siguiente zona del Transformer que puede reelaborarse.
\item La arqueología de algoritmos permite volver a validar, con hardware nuevo y a nueva escala, mecanismos que quedaron olvidados en artículos antiguos.
\item La afinidad con el hardware es el requisito mínimo de la investigación en arquitectura: debe poder paralelizarse, expresarse como multiplicaciones de matrices, implementarse en kernels y escalar (scale).
\item Una de las ventajas de DeepSeek es su fuerte conciencia de la coordinación entre algoritmos e infra/hardware.
\item Los investigadores jóvenes que quieran trabajar en arquitectura necesitan estar cerca de un frontier lab y de un entorno con capacidad de cómputo; de lo contrario, les será difícil validar el comportamiento a escala (scale).
\end{enumerate}

\subsection{Preguntas abiertas}

Las preguntas abiertas se dejan para el final porque la ruta de Attention aún no ha convergido. Las decisiones de Kimi, DeepSeek y MiniMax no son definitivas, sino apuestas provisionales de distintas empresas bajo distintas restricciones.

\begin{itemize}
\item ¿Podrá una Hybrid Attention como Kimi Linear seguir acercándose al desempeño de Full Attention a mayor escala?
\item ¿Podrá Sparse Attention sustituir de verdad a las capas globales sin perder la capacidad de recuperación de largo alcance?
\item ¿La proporción hybrid de tres a uno se consolidará como práctica de largo plazo o será una solución de transición?
\item ¿Se prolongará la vigencia de Full Attention gracias a que el hardware siga optimizándose?
\item Después de MoE, ¿un avance en Attention se convertirá en el siguiente consenso de arquitectura?
\item ¿Cómo puede la investigación en arquitectura construir un puente de validación confiable entre la pequeña escala de código abierto y la gran escala de código cerrado?
\end{itemize}

\subsection{Lecturas adicionales}

\begin{itemize}
\item Kimi Linear: An Expressive, Efficient Attention Architecture.
\item Informes técnicos sobre DeepSeek Sparse Attention / Native Sparse Attention.
\item Serie FlashAttention: para entender la coordinación entre algoritmos y hardware.
\item Materiales sobre arquitecturas como MoE, RoPE, NoPE, Mamba, DeltaNet y Gated Linear Attention.
\item Blogs relacionados de Hazy Research: modelado eficiente de secuencias y arquitecturas afines al hardware.
\end{itemize}

\begin{importantbox}{Valoración final}
Lo que más vale la pena conservar de EP119 es una actitud de investigación: no preguntar solo «qué attention está más de moda», sino qué cuello de botella resuelve, qué capacidad expresiva sacrifica, cómo se paraleliza en el hardware, si supera el scaling ladder y si puede rendir de forma estable en modelos frontier reales.
\end{importantbox}

\end{document}
